% part: intuitionistic-logic
% chapter: soundness-completeness
% section: canonical-model

\documentclass[../../../include/open-logic-section]{subfiles}

\begin{document}

\olfileid{int}{sc}{mod}

\olsection{النموذج القانوني}

نختار لعوالم هذا النموذج متتاليات منتهية~$\sigma$ من الأعداد الطبيعية،
بحيث~$\sigma \in \Nat^*$. وتعريف~$\Nat^*$ استقرائي على الوجه الآتي:
\begin{enumerate}
\item $\emptyseq \in \Nat^*$.
\item متى كان~$\sigma \in \Nat^*$ و$n \in \Nat$، كان
  $\sigma.n \in \Nat^*$؛ ونعني بـ~$\sigma.n$
  $\sigma \concat \tuple{n}$، وبـ~$\sigma \concat \sigma'$ وصل
  $\sigma$ بـ~$\sigma'$.
\item وليس في~$\Nat^*$ شيء وراء ما يقتضيه هذان الشرطان.
\end{enumerate}
ولهذا جاز أن تكون~$\Nat^*$ أساسًا لتعريفات استقرائية.

ولتكن~$\tuple{!B_1, !C_1}$، $\tuple{!B_2, !C_2}$، \dots، تعدادًا
لأزواج ال!!{formula}s كلها. إذا أُعطيت مجموعة
ال!!{formula}s~$\Delta$، عرّفنا~$\Delta(\sigma)$ بالاستقراء التالي:
\begin{enumerate}
\item $\Delta(\emptyseq) = \Delta$.
\item $\Delta(\sigma.n) = {}$
  \[
  \begin{cases}
    (\Delta(\sigma) \cup \{!B_n\})^* &
    \text{إذا كان $\Delta(\sigma) \cup \{!B_n\} \Proves/ !C_n$} \\
    \Delta(\sigma) & \text{خلاف ذلك}
  \end{cases}
  \]
\end{enumerate}
والمراد بـ~$(\Delta(\sigma) \cup \{!B_n\})^*$ مجموعة
ال!!{formula}s الأولية التي نحصل عليها من
\olref[lin]{lem:lindenbaum} حين نطبقها على
$\Delta(\sigma) \cup \{!B_n\}$ وعلى ال!!{formula}~$!C_n$. ومن مقتضى
هذا الاختيار أنه متى كان~$\Delta(\sigma) \cup \{!B_n\} \Proves/ !C_n$
ثبت~$\Delta(\sigma.n) \Proves !B_n$ وبقي
$\Delta(\sigma.n) \Proves/ !C_n$. كذلك يتحقق
$\Delta(\sigma) \subseteq \Delta(\sigma.n)$ لكل~$n$. فإذا كانت
$\Delta$ أولية، لزم أن تكون~$\Delta(\sigma)$ أولية لكل~$\sigma$.

\begin{defn}\ollabel{defn:canonical-model}
  لتكن~$\Delta$ أولية. نسمّي \emph{النموذج القانوني}
  $\mModel{M(\Delta)}$ للمجموعة~$\Delta$ النموذج الذي أجزاؤه كما يأتي:
  \begin{enumerate}
  \item $W = \Nat^*$، أي مجموعة المتتاليات المنتهية من الأعداد الطبيعية.
  \item $R$ ترتيب جزئي يحدّه التكافؤ الآتي:~$R\sigma\sigma'$ إذا وفقط إذا
    كانت~$\sigma$ قطعة ابتدائية من~$\sigma'$؛ ومعنى ذلك أن
    $\sigma'=\sigma\concat\sigma''$ لمتتالية~$\sigma''$ ما.
  \item $V(p) = \Setabs{\sigma}{p \in \Delta(\sigma)}$.
  \end{enumerate}
\end{defn}

والتحقق من كون~$R$ ترتيبًا جزئيًا يسير. وأما شرط الرتابة في~$V$
فحاصل أيضًا: إذ من~$\Delta(\sigma)\subseteq\Delta(\sigma.n)$ يلزم بالاستقراء
$\Delta(\sigma)\subseteq\Delta(\sigma')$ كلما تحقق
$R\sigma\sigma'$.

\end{document}
